\subsection{Run 82: the observable family is a one-parameter arithmetic
curve}\label{sub:r82}

The family of compressed observables studied throughout this programme,
\[
  V_n \;=\; P_+\bigl(T(n)+T(-n)\bigr)P_+ ,
  \qquad n\in(\ZZ/N)^2\setminus\{0\},
\]
carries an exact rigidity that reduces it from $(N^2-1)/K$ nominal members
(the $A^{\mathsf T}$-orbits) to at most $N$.  Let $T\le SL_2(\ZZ/N)$ be the
centraliser of $A$ -- the group whose Weil operators are the \emph{Hecke
operators} of Kurlberg and Rudnick \cite{KR2000} -- and let
\[
  Q(n)\;=\;n_1^2-n_1n_2-n_2^2 \pmod N
\]
be its invariant quadratic form (discriminant $t^2-4=5$ at $t=3$).

\begin{theorem}[torus rigidity]\label{thm:torus}
Let $N$ be an odd prime with $\gcd(N,t^2-4)=1$.  If $Q(n)=Q(n')$ then
$V_{n'}=U_gV_nU_g^{*}$ for an explicit $g=aI+bA\in T$; in particular
$\operatorname{spec}(V_n)=\operatorname{spec}(V_{n'})$.  Consequently the
family has at most $N$ distinct spectra, one for each value of $Q$; each
nonzero level set of $Q$ is a single $T$-orbit containing exactly $|T|/K=m$
distinct $A^{\mathsf T}$-orbits, and
\[
  (N\mp1)\,m \;+\; \frac{|\{Q=0\}\setminus\{0\}|}{K}\;=\;\frac{N^2-1}{K}.
\]
\end{theorem}

\begin{proof}
Identify $(\ZZ/N)^2$ with $R_A=(\ZZ/N)[x]/(x^2-tx+1)$ as a free rank-one
$R_A$-module.  Then $Q$ is a scalar multiple of the norm form of $R_A$ and $T$
is its group of norm-one units, so $T$ acts simply transitively on each level
set $\{Q=q\}$, $q\neq0$; in particular $n'=gn$ for some $g\in T$.  Since $g$
commutes with $A$, the Weil operator $U_g$ commutes with $U_A$ up to a scalar
and therefore preserves every $U_A$-eigenspace; since $-I$ is central, $U_g$
commutes with the parity operator and preserves each parity block.  The Weil
intertwining property $U_gT(n)U_g^{*}=T(gn)$ then gives
$U_gV_nU_g^{*}=P_+(T(gn)+T(-gn))P_+=V_{n'}$.
\end{proof}

Theorem \ref{thm:torus} is elementary given \cite{KR2000}, but it settles a
question this programme had left open, and it constrains the ensemble route
sharply.

\begin{corollary}\label{cor:curve}
The $n$-ensemble $\{V_n\}$ is a deterministic one-parameter curve
$q\mapsto V_q$, each member repeated $m$ times, and every average over
$n\in(\ZZ/N)^2$ equals the corresponding average over $q\in\FF_N^{*}$.
Hence \textup{(i)} the ensemble has one intrinsic fluctuation direction, so a
single dominant coherent mode and positively correlated bulk window counts are
automatic; \textup{(ii)} no i.i.d.\ $\GOE$ statement can hold for $\{V_n\}$;
\textup{(iii)} averaging over $n$ buys concentration in one arithmetic scalar
only and therefore cannot substitute for a per-system local law.
\end{corollary}

Part \textup{(iii)} closes, for this family, the last route in which an
ensemble average might have replaced a local law; together with the
moment-route obstruction it leaves the single-matrix problem: for fixed $q$,
prove $\GOE$ local statistics for $V_q$.  In the Hecke basis $\{f_j\}_{j\in
\ZZ/m}$ (the $T$-character basis) equivariance forces
$\langle f_j|T(n)|f_k\rangle=c_{jk}(q)\,\zeta^{(j-k)s(n)}$ with
$\zeta=e^{2\pi i/m}$ and $s(n)$ the discrete logarithm of the torus element
carrying a base point to $n$; so
$V_n=D_{s(n)}(C_q+C_q^{*})D_{s(n)}^{*}$ with $D_s$ diagonal, and the entire
arithmetic content is one matrix $C_q=(c_{jk}(q))$ of explicit exponential
sums per level set.

\emph{The modulation, measured exhaustively.}  With $q$ as the parameter the
curve is computable over \emph{all} level sets.  At $(N,K)=(3571,17)$,
$m_+=105$, over all $3570$ values of $q$: $\operatorname{Tr}(V_q^2)/m_+$ has
mean $0.117600$ against the arithmetic value $N/(K^2m_+)=0.117658$ and
relative spread $1.22\%$, while
$\mu_4(q)\in[1.2246,2.8573]$ with mean $1.9694$ and standard deviation
$0.2350$ (the semicircle value is $2$; the arcsine value $3/2$ is excluded).
The corresponding matched $\GOE$ reference has $\operatorname{std}(\mu_4)
=0.027$, so the macroscopic shape of $V_q$ does not self-average.  Writing
$A^j=u_jI+v_jA$, so that $v_j=1,3,8,21,55,144,377,987,\dots$ (the even-index
Fibonacci numbers at $t=3$), the additive Fourier transform in $q$ of the raw
central traces is supported on those coefficients: the $16$ frequencies
$\pm v_j$, $1\le j\le 8$, carry $90\%$ of the variance of
$\operatorname{Tr}(V_q^4)$ and $\operatorname{Tr}(V_q^6)$, and a
least-squares fit on exactly those $(K-1)/2$ frequency pairs gives
$R^2=0.9426$, $0.9276$ and, for $\mu_4(q)$, $R^2=0.98655$.  Enlarging the
frequency set to $71$ frequencies moves $R^2$ only to $0.9472$/$0.9872$: the
comb is saturated.  The shape is that of a weighted sum of $(K-1)/2$
oscillations, matching the form conjectured by Kurlberg and Rudnick for the
matrix elements of the desymmetrised cat map \cite{KR2003}; the coefficients
are fitted here, not derived.

\emph{Where the excess lives.}  In the basis diagonalising linear eigenvalue
statistics, $L_k=\sum_iT_k(\lambda_i/\sqrt2)$, the variance excess over the
matched $\GOE$ reference is confined to \emph{even} modes: the ratios at
$k=4,6,8,10$ are $81,71,78,100$ at $m_+=105$ and $375,306,302,352$ at
$m_+=246$, against $1.0,1.6,2.8,5.8$ and $0.5,2.9,5.2,10.8$ at
$k=3,5,7,9$; a Haar control on the same operators gives $0.3$--$0.6$
throughout.  Removing the component of the bulk count vectors explained by
$k\le8$ drops the count-variance ratio from $1.89$ to $1.12$ ($m_+=105$) and
from $3.82$ to $1.46$ ($m_+=246$).  The excess is not an edge effect: the
largest normalised eigenvalue sits at $1.35$ bulk edges on average against
$\GOE$'s $1.39$, and trimming the eight most deviant eigenvalues leaves the
$\mu_4$ variance ratio at $52$ (from $69$).  This is exactly the signature of
a one-parameter family of symmetric macroscopic densities.
